\scanpage{21}\mspage{17}
% @ top, end of lemma | text | E(u), ∃! base
$E(u)$, $\exists!$ base $b\subset C$ telle que $b\cap E(u)=\{\unsure{\zeta|\xi}_1,\zeta_2\}$ et (si $b'$ est la base opposée) on a $b'\cap E(u)=\{\sigma\zeta_1,\sigma\zeta_{\unsure{1|2}}\}$, \deleted{\unsure{donc}} où $\sigma$ est l'autom can. d'ordre 2 de $E(u)$.

% @ above proof, NB | text | NB On aura alors
$\uparrow$ \uline{NB} On aura alors $\sigma\zeta_1=\zeta'_2$, $\sigma\zeta_2=\zeta'_1$. Alors $b\cap E(u)=\{\zeta_1,\zeta_2\}$, $b'\cap E(u)=\{\sigma\zeta_1,\sigma\zeta_2\}$.

% @ upper third, demonstration | text | Dém Soit une base
\uline{Dém} Soit une base $\{\zeta_1,\zeta_2,-\zeta_6\}$ satisfaisant à \unsure{l'énoncé}.
\margin{\unsure{principe|rappeler} plus simple plus bas}% slanted
On aura alors \deleted{\unsure{tel que}} $u$ est lié à $\zeta_1,\zeta_2,\zeta'_1,\zeta'_{\unsure{2}}$, donc $u=\zeta_{12}$, et $b\cap E(u)=\{\zeta_1,\zeta_2\}$, $b'\cap E(u)=\{\zeta'_1,\zeta'_2\}$. Pour l'existence \inserted{d'une base $b$ satisfaisant les conditions envisagées}, on note que (comme le groupe $W_u$ est transitif sur l'ens des \deleted{couples} \inserted{paires} $\{\zeta_1,\zeta_2\}$ non liées de $E(u)$ et $W=W_{E_6}$ transitif sur l'ens des $u\in C$) on peut supposer, p.r. à une base $(\bar\zeta_1,-\bar\zeta_6)$ déjà donnée, que $u=\bar\zeta_{12}$, $\zeta_1=\bar\zeta_1$, $\zeta_2=\bar\zeta_2$.

% @ centre and crossed region | text | Pour l'unicité, on note que
Pour l'unicité, on note que \deleted{\unsure{alors que} \unsure{les} $\bar\zeta'_i$ \unsure{sont} $(3\le i\le6)$ \unsure{si} $\bar\zeta_{\unsure{3}}$ est \unsure{un autre él.|un autre}  lié à $\zeta_2$ et $\unsure{\sigma\zeta_1|\zeta_1|\sigma\xi_1}=\zeta'_2$ et $\sigma\zeta_2=\zeta'_1$, et un él. de $C$ non lié à $\zeta_1,\zeta_2$ distinct de $\zeta_1$, donc est \unsure{dans|au} $\complement E(u)$, et de \unsure{même pour les} d'après ce qu'on a vu, on doit avoir que les $\bar\zeta_i,\bar\zeta'_i$ $(3\le i\le6)$ sont $\in\complement E(u)$, donc qu'ils \inserted{ceux qui} ne sont pas des $\zeta_j,\zeta'_j$,} alors que $\bar\zeta_3$ doit être non lié à $\zeta_1,\zeta_2$ et lié à $\zeta'_1,\zeta'_2$ donc ce n'est pas un $\zeta'$, ni un $\zeta_{ij}$ (car ce serait un $\zeta_{1j}$ ou $\zeta_{2j}$ pour être lié à $\zeta'_1,\zeta'_2$, mais il le serait aussi à $\zeta_1,\zeta_2$, absurde), de $=$ pour les $\bar\zeta_i$ $(3\le i\le6)$. cqfd.

% @ left of cancelled lines | diagram | two triangles sharing lower vertex, ζ₁, σζ₁, ζ₂, σζ₂
\diagram{\begin{tikzcd}[ampersand replacement=\&,execute at end picture={
\draw (0,0)--(.28,1.2) node[above left] {$\zeta_1$} --(.70,1.14) node[above right] {$\sigma\zeta_1$} --cycle;
\draw (0,0)--(1.43,.92) node[above] {$\zeta_2$} --(1.65,.54) node[right] {$\sigma\zeta_2$} --cycle;
}]\vphantom{A}\end{tikzcd}}

% @ foot, NB | text | Positions relatives
\uline{NB} Positions relatives \inserted{(à conjugaison près)} d'un $u\in C$ et d'une base $b\subset C$.

% @ foot, first case | text | (i) u ∈ b
(i) $u\in b$ \quad Pour $u\inserted{=\zeta_1}$ fixé, il y a \unsure{18|16} telles bases, corresp. \unsure{aux} au choix d'un $\zeta'_1\in C-E(u)$
\margin{$\left.\begin{aligned}r_b\unsure{(u)|{}\cdot u}&=1\\E(u)\cap b&=\varnothing\\\operatorname{card} E(u)\cap b'&=5\end{aligned}\right\}$}

% @ foot, second case | text | (ii) u ∈ b′
\deleted{\unsure{(ii)}} (ii) $u\in b'$ (base opposée) Pour $u\inserted{=\zeta'_1}$ fixé, 16 telles bases, correspondant au choix d'un $\zeta_1\in C-E(u)$
\margin{$\left.\begin{aligned}r_b\cdot u&=-1\\\operatorname{card} E(u)\cap b^{\unsure{+|*|{}}}&=5\\E(u)\cap b'&=\varnothing\end{aligned}\right\}$}

% @ foot, third case | text | (iii) u ∉ b ∪ b′
(iii) $u\notin b\cup b'$ \quad Pour $u$ fixé, il y a 40 telles bases ($r_b$ \uline{orthogonal} à $u$)
\margin{$\left.\begin{aligned}r_b\cdot u&=0\\\operatorname{card} E(u)\cap b&=2\\\operatorname{card} E(u)\cap b'&=2\end{aligned}\right\}$}

% @ foot right, total | display | 16+16+40
\[
\overline{\overline{16+16+40}}=72
\]
% archive sheet number: 20
